% !TEX root = ../main.tex
\section{Kết luận và Hướng phát triển}

\subsection{Kết luận}
Đồ án này đã tập trung tìm hiểu vấn đề "Khởi động lạnh" (Cold-start), một khó khăn kỹ thuật làm giảm hiệu năng của mô hình Điện toán phi máy chủ (Serverless Computing). Thông qua việc khảo sát nền tảng Knative, chúng tôi nhận thấy cơ chế mở rộng phản ứng (Reactive Autoscaling) thường không phản hồi kịp thời khi lưu lượng truy cập mạng tăng đột biến.

Để cải thiện tình trạng này, một hệ thống Lập lịch Dự đoán (Predictive Scheduling) đã được xây dựng. Thay vì đợi tải tăng, hệ thống sử dụng mạng nơ-ron LSTM để tính toán trước lượng request, từ đó điều khiển Kubernetes tạo sẵn Pod thông qua một Custom Controller. Ngoài ra, một thuật toán làm mịn (Smoothing Algorithm) đơn giản cũng được tích hợp để tránh việc tạo/xóa Pod liên tục làm quá tải hệ thống.

Kết quả thử nghiệm bước đầu cho thấy hệ thống có khả năng dự đoán đúng các đợt tải, giúp giảm thiểu đáng kể số lượng request bị nghẽn do khởi động lạnh. Dù vẫn còn hiện tượng lãng phí một chút tài nguyên do dự đoán dư, nhưng nhìn chung giải pháp đã phần nào cân bằng được giữa chi phí và tốc độ phản hồi.

\subsection{Hướng phát triển tương lai}
Vì điều kiện và thời gian có hạn, đồ án vẫn còn một số điểm cần khắc phục và cải thiện trong thời gian tới:

\begin{enumerate}
    \item \textbf{Đưa lên môi trường thực tế:} Đồ án hiện chỉ chạy thử nghiệm trên máy tính cá nhân (Local Cluster) với môi trường giả lập. Bước tiếp theo là triển khai hệ thống lên các nền tảng đám mây thực tế (AWS hoặc Google Cloud) với quy mô nhiều máy chủ (Multi-node) để đo đạc giới hạn chịu tải.
    \item \textbf{Cải tiến mô hình học máy:} Mô hình LSTM hiện tại chỉ đang xử lý dữ liệu đầu vào là số lượng request (RPS). Nếu có thể áp dụng các mô hình mạnh hơn như Transformer để phân tích thêm nhiều loại số liệu khác (CPU, RAM, độ trễ mạng), kết quả dự báo có thể sẽ chính xác hơn rất nhiều.
    \item \textbf{Quản lý tài nguyên tự động hơn:} Thuật toán làm mịn số lượng Pod hiện đang được code tay (Hard-coded heuristics) dựa trên quan sát thực nghiệm. Về lâu dài, có thể áp dụng Học tăng cường (Reinforcement Learning) để hệ thống tự học cách cân bằng giữa việc tiết kiệm tài nguyên và đảm bảo tốc độ.
\end{enumerate}

Nhìn chung, việc đưa Machine Learning vào hỗ trợ vận hành Serverless là một hướng đi rất khả thi. Đồ án này là bước đệm cơ bản để tiếp tục phát triển các hệ thống quản lý tài nguyên thông minh và linh hoạt hơn.